% !PRINTED 542
\weberpage{542}
Sind dann $k$, $k_0$, $k'$ die Klassen, zu denen $\psi$, $\psi_0$, $\psi'$ gehören, so ist
\[
\chi(\delta, k') = \chi(\delta, k_0)\,\chi(\delta, k),
\]
also
\weq{5}{\begin{aligned}
\chi(\delta, k) &= (\delta, A), \\
\chi(\delta, k') &= (\delta, 2)(\delta, A).
\end{aligned}}

Ist $\omega_1$ eine Wurzel der Klasse $k$, so ist $\omega_2$ Wurzel der entgegengesetzten Klasse $k^{-1}$, in diesen beiden Klassen sind aber die Charaktere $\chi(\delta, k)$ und $\chi(\delta, k^{-1})$ dieselben. Multiplizieren wir daher die Formel \S\,141, (30) mit $\chi(\delta, k)$ und bilden die Summe über alle Klassen $k$, so folgt:
\weq{6}{\overset{k}{\Sigma}\,\chi(\delta,k)\Bigl(\overset{x,y}{\Sigma}\frac{1}{\psi^s} - \overset{x,y}{\Sigma}\frac{1}{\psi'^s}\Bigr) = \frac{4\pi}{\sqrt{m}}\overset{k}{\Sigma}\,\chi(\delta,k)\log\frac{f(\omega)}{\sqrt[4]{2}},}
und wegen (5) kann man für die linke Seite schreiben:
% [sic] The printed fraction chi(delta,k)/1 is retained; its denominator is redundant, not an algebraic change.
\[
[1-(\delta,2)]\overset{k}{\Sigma}\frac{\chi(\delta,k)}{1}\overset{x,y}{\Sigma}\frac{1}{\psi^s}.
\]
Demnach haben wir:
\weq{7}{\begin{aligned}
(\delta,2) = +1\colon &\quad \frac{2\pi}{\sqrt{m}}\overset{k}{\Sigma}\,\chi(\delta,k)\log\frac{f(\omega)}{\sqrt[4]{2}} = 0, \\
(\delta,2) = -1\colon &\quad \frac{2\pi}{\sqrt{m}}\overset{k}{\Sigma}\,\chi(\delta,k)\log\frac{f(\omega)}{\sqrt[4]{2}} = \overset{k}{\Sigma}\,\chi(\delta,k)\overset{x,y}{\Sigma}\frac{1}{\psi^s},
\end{aligned}}
wobei rechts der Grenzwert für $s=1$ zu nehmen ist.

In \S\,113 haben wir die Formel bewiesen:
\weq{8}{\overset{k}{\Sigma}\,\chi(\delta,k)\overset{x,y}{\Sigma}\frac{1}{\psi^s} = \Sigma\frac{(\delta,n)}{n^s}\Sigma\frac{(\delta',n)}{n^s}.}

Wenn wir mit $K(\delta)$ die Klassenzahl für die Diskriminante $\delta$ bezeichnen, so ist, wie in \S\,112, (8) bewiesen ist: % faint pre-delta speck is scan/print damage, not a character
\weq{9}{\begin{aligned}
\sum^n \frac{(\delta,n)}{n} &= \frac{\pi}{\sqrt{-\delta}}\,K(\delta) \qquad \delta<0, \\
\sum^n \frac{(\delta,n)}{n} &= \frac{\log\varepsilon}{\sqrt{\delta}}\,K(\delta) \qquad \delta>0,
\end{aligned}}
worin die Quadratwurzeln positiv zu nehmen sind, und
\[
\varepsilon = \frac{T+U\sqrt{\delta}}{2}
\]
die fundamentale positive Einheit des quadratischen Körpers mit der Grundzahl $\delta$ ist. Für die beiden Ausnahmefälle $\delta=-3$, $\delta=-4$ gelten diese Formeln, wenn man unter $K$ nicht die
%CONT>>
